\chapter{Identification, statistical efficiency and machine learning}
\label{r3:learning}
\section{Four different learning problems}
Suppose Chapter~\ref{r3:graphs} has justified adjustment by $X$. The
observed data for an independent household are $O=(X,A,Y)$, with binary
$A$, propensity $e(x)=P(A=1\mid X=x)$ and response surfaces
$m_a(x)=\E[Y\mid A=a,X=x]$. The scientific target is the population
average $\theta=\E[m_1(X)-m_0(X)]$, or the conditional contrast
$\tau(x)=m_1(x)-m_0(x)$. These are causal targets only under the
identification assumptions. The same formulas remain statistical
association functionals when those assumptions fail.

Machine learning can estimate nuisance functions $(e,m_0,m_1)$, estimate
the effect function $\tau$, optimize a policy using the estimated effects,
or propose a causal representation or graph. Each has its own loss and
generalization question. An outcome predictor minimizes error for observed
$Y$ under the historical action distribution. A policy learner values
outcomes under changed actions. The two losses are not interchangeable,
and a random train/test split cannot by itself bridge them.

\section{Deriving an influence function for the identified mean}
\label{r3:influence}
Work initially in an unrestricted dominated observed-data model with
$\epsilon<e(X)<1-\epsilon$ and $\E[Y^2]<\infty$. A regular parametric
submodel $P_t$ through $P=P_0$ has mean-zero score
$S(O)=\left.\partial_t\log p_t(O)\right|_{t=0}$. Factor the observed
law as $p(x)p(a\mid x)p(y\mid a,x)$ and decompose
\[
 S=S_X(X)+S_A(A,X)+S_Y(Y,A,X),
\]
where $\E[S_X]=0$, $\E[S_A\mid X]=0$ and
$\E[S_Y\mid A,X]=0$. These are orthogonal score components in
$L_2(P)$, associated with the marginal covariate law, action mechanism
and outcome mechanism.

Differentiating $\theta(P_t)=\int(m_{1,t}-m_{0,t})\,dP_{X,t}$ gives
\begin{align}
 \dot\theta
 &=\E[(m_1-m_0-\theta)S_X]
   +\E[\dot m_1(X)-\dot m_0(X)],\nonumber\\
 \dot m_a(x)
 &=\E[(Y-m_a(x))S_Y\mid A=a,X=x].
 \label{r3:path-derivative}
\end{align}
The conditional mean derivative follows by differentiating its normalized
conditional density; subtracting $m_a$ uses the zero conditional mean
of $S_Y$. The action score does not enter the functional's derivative.
Moving the conditional outcome terms into an expectation over $O$ yields
\begin{equation}
 \phi(O)=m_1(X)-m_0(X)-\theta
       +\frac{A}{e(X)}\{Y-m_1(X)\}
       -\frac{1-A}{1-e(X)}\{Y-m_0(X)\}.
 \label{r3:eif}
\end{equation}
Indeed $\E[\phi]=0$ and $\dot\theta=\E[\phi S]$ for every such score:
the residual terms supply $S_Y$, the centered contrast supplies $S_X$,
and conditioning shows $\E[\phi S_A]=0$. In the unrestricted
nonparametric model the tangent space is $L_2^0(P)$, so this gradient is
the canonical gradient. Its variance is the efficiency bound per
independent observation. This supplies the construction behind the
orthogonal ATE score used by
\citet[Section 5.1, Equation (5.3)]{dml2018}.

Two cautions follow from the derivation. The statistical model and sampling
unit determine the bound; it is not automatically the bound for a panel
with common shocks. Also, pointwise $\tau(x)$ at a continuously distributed
$x$ generally lacks this regular root-$n$ representation without additional
structure or smoothing. A scalar ATE theorem cannot be used as a CATE
confidence-band theorem.

\section{Orthogonality, double robustness and the exact remainder}
\label{r3:ml-remainder}
Write $\Phi=\phi+\theta$ and replace the nuisances by estimates from
a separate training sample. Conditional on that training sample,
\begin{equation}
 P\widehat\Phi-\theta
 =\E\!\left[(\widehat e-e)
       \left\{\frac{\widehat m_1-m_1}{\widehat e}
             +\frac{\widehat m_0-m_0}{1-\widehat e}\right\}\right].
 \label{r3:exact-remainder}
\end{equation}
For the treated term, conditioning on $X$ gives
$\widehat m_1+e(m_1-\widehat m_1)/\widehat e-m_1
=(\widehat e-e)(\widehat m_1-m_1)/\widehat e$.
The untreated term is analogous with the sign reversed by the ATE
contrast. This proves the displayed identity.

The bias is zero if the propensity is correct, or if both outcome
regressions are correct. That is double robustness at the population
moment level. At the joint truth the derivative in every nuisance
direction is also zero: the error is second order. That is Neyman
orthogonality. The concepts are related but not identical: an orthogonal
score need not have the union-model robustness expressed by this exact
identity.

If fitted propensities lie in $[\epsilon,1-\epsilon]$,
Cauchy--Schwarz gives
\[
 |P\widehat\Phi-\theta|
 \le\epsilon^{-1}\|\widehat e-e\|_2
       \bigl(\|\widehat m_1-m_1\|_2+
             \|\widehat m_0-m_0\|_2\bigr).
\]
Both nuisances having $o_p(n^{-1/4})$ error is one sufficient route to
$o_p(n^{-1/2})$ bias; it is not necessary that both achieve the same
rate. A known randomized propensity makes this particular bias exactly
zero for a fixed fitted outcome function. Finite variance and the
sampling conditions still matter.

For fixed $K$ folds, fit nuisances outside each fold and average
$\widehat\Phi_i$ inside it. On each evaluation fold,
\[
 \widehat\theta-\theta
 =(P_n-P)\Phi+P(\widehat\Phi-\Phi)
                  +(P_n-P)(\widehat\Phi-\Phi).
\]
The first term is the empirical average of $\phi$. The second is the
remainder above. Conditional independence of training and evaluation
data, together with $L_2$ consistency of the score, controls the third
term at $o_p(n^{-1/2})$. Appropriate moments, positive finite limiting
variance and the nuisance product rate give
\[
 \sqrt n(\widehat\theta-\theta)\ \Longrightarrow\
 N(0,\E[\phi^2]).
\]
The assumptions are the substantive content of the DML guarantee
\citep[Assumption 5.1 and Theorem 5.1]{dml2018}; cross-fitting alone
is not a theorem.

Repeated records from one customer remain one dependent trajectory.
Independent-trajectory splitting holds out the entire customer. Common
calendar shocks can couple even different customers, requiring a different
sampling theorem and covariance calculation. For a single mixing time
series, separated blocks may control training/evaluation dependence under
quantified mixing and gap conditions. A shuffled row split supplies no
such argument. The proposal must attach fold construction and the
appropriate CLT to each estimator, rather than expose a universal
``cross-fit'' switch.

\section{Learning heterogeneous effects: compare the losses}
\label{r3:heterogeneity}
The statistical CATE is a difference of response surfaces, but that does
not determine the best finite-sample estimation strategy. Several methods
can target the same $\tau$ while exploiting different structure.

\subsection{Response models and imputed contrasts}
A T-learner fits $m_0,m_1$ separately. An S-learner fits a common
regression $m(a,x)$ and evaluates its difference. Sharing parameters may
help when arms are similar, but regularization can suppress action
interactions even when the factual prediction loss barely changes. A
T-learner can waste information when the arms share a complicated
baseline but the effect is simple.

An X-learner first fits both response surfaces and constructs
$D_i^1=Y_i-\widehat m_0(X_i)$ among treated units and
$D_i^0=\widehat m_1(X_i)-Y_i$ among controls. It then regresses those
imputed contrasts on $X$ within their respective arms and combines the
resulting effect regressions with a chosen weight $w(x)$:
\[
 \widehat\tau_X(x)=w(x)\widehat\tau_0(x)
                         +(1-w(x))\widehat\tau_1(x).
\]
Imbalance can favor this construction if the larger arm estimates one
response surface well and the effect function is easier to learn. The
weight is a statistical choice, not proof of exchangeability
\citep[Section on the X-learner]{kunzel2019}.

\subsection{Residual learning and its population metric}
Let $m(X)=\E[Y\mid X]=m_0(X)+e(X)\tau(X)$. Subtraction gives
\begin{equation}
 Y-m(X)=(A-e(X))\tau(X)+\varepsilon,
 \qquad \E[\varepsilon\mid A,X]=0.
 \label{r3:robinson}
\end{equation}
The R-learner minimizes a cross-fitted, regularized version of
\[
 L(f)=\E[\{Y-m(X)-(A-e(X))f(X)\}^2].
\]
At the true nuisances the cross term vanishes, so
\begin{equation}
 L(f)-L(\tau)=\E[e(X)(1-e(X))\{f(X)-\tau(X)\}^2].
 \label{r3:r-risk}
\end{equation}
The loss concentrates information where both arms occur. Its curvature
vanishes where treatment is deterministic. The quasi-oracle analysis of
\citet[Sections 2 and 5]{rlearner2020} makes nuisance-rate and
function-class conditions explicit; Equation~\eqref{r3:r-risk} alone
does not establish a rate for an arbitrary neural learner.

If $f$ is constrained to a constant, the target becomes
\[
 \theta_R=\frac{\E[e(X)(1-e(X))\tau(X)]}
                 {\E[e(X)(1-e(X))]},
\]
which need not equal the ATE. For two equally common groups with effects
$0$ and $10$ and propensities $0.5$ and $0.01$, respectively, the ATE
is $5$ but the residual-loss constant is approximately $0.381$.
Neither estimate is a numerical error if it answers its specified
projection target. This target change must be visible in a product
comparing a constant partially linear model with heterogeneous effects.

\subsection{Doubly robust pseudo-outcomes}
The DR-learner regresses $\widehat\Phi(O)$ from
Equation~\eqref{r3:eif}, without subtracting $\theta$, on $X$.
At true nuisances, $\E[\Phi\mid X]=\tau(X)$. With estimated nuisances,
the conditional bias is exactly the integrand of
Equation~\eqref{r3:exact-remainder}. The second-stage regression can
adapt to a simpler effect function, but its error includes both that
regression's stochastic error and how it smooths the nuisance-product
bias. A product of global $L_2$ errors sufficient for an ATE does not
by itself guarantee a pointwise CATE rate.

\citet[Sections 3--4]{drlearner2023} develops stability and oracle-error
conditions for regression with estimated pseudo-outcomes. That formulation
clarifies why using an estimated outcome is not a free replacement for
observing the individual treatment effect. It also distinguishes the
second-order error of an orthogonal loss from the double-robust product
structure of this pseudo-outcome.

\subsection{Why a local effect has a different rate from an average}
\label{r3:cate-rates}
To expose the distinction, take an independent nuisance-training sample
and a linear second-stage smoother
$\widehat\tau(x)=\sum_iw_i(x)\widehat\Phi_i$, with weights determined
by the second sample's covariates. Its oracle analogue uses the true
pseudo-outcomes, $\widetilde\tau(x)=\sum_iw_i(x)\Phi_i$.
Conditional on the training data and evaluation covariates,
\[
 \widehat\tau(x)-\widetilde\tau(x)
 =\sum_iw_i(x)b(X_i)+\sum_iw_i(x)\xi_i,
 \quad \E[\xi_i\mid X_i,\text{training}]=0,
\]
where $b$ is the conditional nuisance-product bias above. If
$\sum_i|w_i(x)|=O_P(1)$,
$\sum_iw_i(x)^2=O_P((nh^d)^{-1})$, and the local conditional
second moment of $\widehat\Phi-\Phi$ is bounded by $\delta_n^2\to0$,
the last term is $O_P\{\delta_n(nh^d)^{-1/2}\}$ by conditional
variance and Chebyshev's inequality. Under bounded overlap the first
is controlled by a local weighted product of nuisance errors. This
project argument gives a concrete sufficient stability condition;
\citet[Theorem 2]{drlearner2023} formulates a more general one.

Suppose the CATE is $\gamma$-smooth in $d$ covariates and a regular
local polynomial design gives oracle bias $O(h^\gamma)$ and
standard deviation $O((nh^d)^{-1/2})$. Balancing them uses
$h\asymp n^{-1/(2\gamma+d)}$ and yields oracle error
$n^{-\gamma/(2\gamma+d)}$. If the propensity and outcome regressions
are respectively $\alpha$- and $\beta$-smooth, and attain their
corresponding rates in the \emph{local weighted norms} needed by the
smoother, the comparison becomes
\begin{equation}
 \widehat\tau(x)-\tau(x)
 =O_P\left(n^{-\gamma/(2\gamma+d)}
       +n^{-\alpha/(2\alpha+d)-\beta/(2\beta+d)}\right).
 \label{r3:cate-rate}
\end{equation}
This is the rate in \citet[Corollary 1]{drlearner2023}, under its
stability, pseudo-outcome consistency and optimal-smoother conditions.
A sufficient strict condition for negligible nuisance bias is
$\alpha/(2\alpha+d)+\beta/(2\beta+d)>\gamma/(2\gamma+d)$.
The analogous root-$n$ average-effect argument needs the right side
replaced by $1/2$. Conversely, global $L_2$ rates can conceal poor
behavior near one target $x$ and do not establish the local premise.

Oracle equivalence transfers an oracle inferential result only when
that result is itself valid. At a mean-squared-error-optimal bandwidth,
smoothing bias can be of the same order as standard deviation;
inference may require bias correction or undersmoothing. Forest
intervals, kernel intervals and RKHS excess-risk theorems therefore
need distinct guarantees. A score being orthogonal is insufficient
to label every displayed customer-level curve a calibrated causal
confidence band. Simultaneous coverage across many customer strata
is an additional uniform-inference problem.

\subsection{Forests as local moment estimators}
A generalized random forest constructs neighborhood weights
$\alpha_i(x)$ from tree leaves and solves
\[
 \sum_i\alpha_i(x)\psi(O_i;\tau(x),\nu(x))=0,
 \qquad \sum_i\alpha_i(x)=1.
\]
The score can be a local outcome, treatment or instrumental moment. Thus
the same computational architecture can estimate different statistical
objects. Identification still depends on which moment is supplied.
Honesty separates observations used to choose splits from those used
to estimate leaf moments. Subsampling, regular splits, smoothness,
local identification and further conditions enter the forest CLT
\citep[Algorithm 1, Specification 1 and Theorem 5]{grf2018}. Default
forest fitting does not automatically establish all these assumptions,
and local weak instruments remain weak even when the forest predicts
receipt well elsewhere.

\begin{center}\small
\begin{tabularx}{\textwidth}{P{23mm}XX}
\toprule
Construction & Where structure is learned & Consequential limitation\\
\midrule
T/S learners & Factual response surfaces & Small difference can be
obscured by response-model regularization.\\
X-learner & Imputed arm-specific contrasts & Depends on the first-stage
surface and the arm-imbalance/effect structure.\\
R-learner & Residual loss for effect function & Curvature and projection
weights depend on overlap.\\
DR-learner & Conditional mean of orthogonal pseudo-outcome & Requires a
second-stage regression analysis, not just an ATE remainder bound.\\
Generalized forest & Local moment and adaptive neighborhood & Honesty,
smoothness, sampling and local rank govern inference.\\
\bottomrule
\end{tabularx}
\end{center}

\section{Targeting, policy learning and the value of decisions}
\label{r3:targeting}
A one-step ATE estimate averages the augmented expression. A targeted
substitution estimator instead updates a fitted outcome regression so
that its empirical score is close to zero, then averages the updated
contrast. For binary outcomes, one possible fluctuation is
\[
 \operatorname{logit}m_{a,\varepsilon}(x)
 =\operatorname{logit}\widehat m_a(x)+\varepsilon H(a,x),\quad
 H(a,x)=\frac a{\widehat e(x)}-\frac{1-a}{1-\widehat e(x)}.
\]
The binomial likelihood derivative is
$\sum_iH(A_i,X_i)(Y_i-m_{A_i,\varepsilon}(X_i))$.
Solving it and taking
$\widehat\theta=P_n(m_{1,\widehat\varepsilon}-m_{0,\widehat\varepsilon})$
makes the empirical influence-function mean zero for these fits. The
logistic map respects outcome bounds. This calculation explains a
targeting step. Bernoulli cross-entropy remains a valid mean loss for
any outcome in $[0,1]$, even when it is not binary: its conditional
expected derivative vanishes at the conditional mean. A financial
outcome with known fixed bounds can first be rescaled. Clipping an
unbounded monetary outcome, however, changes the estimand and cannot
be justified merely to use a logistic implementation.

To connect the update to inference, let $m^*$ denote the targeted fits
and $\widehat\Phi^*$ their augmented expression. The solved score gives
$P_n\widehat\Phi^*=P_n(m_1^*-m_0^*)=\widehat\theta$.
Consequently the expansion is still the empirical true influence
function plus the exact remainder of
Equation~\eqref{r3:exact-remainder}, now evaluated at $m^*,\widehat e$,
and an empirical-process difference. The product-rate condition must
hold for the updated fits. If initial training uses separate folds but
the fluctuation is fitted on the evaluation observations, the final fits
are not literally independent of those observations. Their additional
finite-dimensional dependence must be controlled, or a specified
cross-validated targeting construction used. The longitudinal
construction and the score/remainder analysis appear in
\citet[Chapter 4, Sections 4.4--4.5]{targeted2018}; its displayed
Donsker condition is one empirical-process route, not an automatic
property of an arbitrary neural learner.

For a fixed binary policy $d(X)$, and a per-action cost $c(X)$ in the
same units as the outcome, incremental value over no treatment is
\[
 V(d)-V(0)=\E[d(X)\{\tau(X)-c(X)\}].
\]
Without capacity constraints the ideal rule is
$d^*(x)=\ind\{\tau(x)>c(x)\}$. A budget constraint
$\E[d(X)b(X)]\le B$ introduces a multiplier and changes the threshold
to $\tau(x)>c(x)+\lambda b(x)$ under the usual separable relaxation.
Optimizing an estimated score within a restricted policy class differs
from estimating a full accurate CATE surface: regret weights errors by
their effect on the decision boundary and by the policy restrictions.
Policy evaluation and policy learning therefore have distinct guarantees
\citep{policy2021}.

If the same data choose the winning policy and produce its ordinary
fixed-policy interval, the interval does not account for selection.
An independent evaluation sample, valid nested splitting, or uniform
policy-class inference is required. A model can rank policies usefully
without estimating every customer effect accurately, but a claim of
positive incremental value still needs uncertainty for the chosen
decision procedure and the appropriate target population.

\section{Representations do not manufacture an adjustment set}
\label{r3:representation}
Suppose $X$ suffices for exchangeability and a learned representation
$B=b(X)$ is proposed as a replacement. A useful sufficient condition is
that it also balances action: $A\perp X\mid B$. Indeed,
\[
 P\{Y(a)\mid A,B\}
 =\int P\{Y(a)\mid x\}\,dP(x\mid A,B)
 =\int P\{Y(a)\mid x\}\,dP(x\mid B).
\]
The first equality uses exchangeability given $X$ and the fact that
$B$ is a function of $X$; the second uses balance. Without that extra
property a lossy representation can destroy the identification route.
The true propensity $e(X)$ balances binary treatment because conditioning
on $e(X)=b$ makes the probability of treatment the same $b$ for every
remaining $X$ stratum. An approximate learned score needs an estimation
and overlap analysis.

In the confounded two-world example of Chapter~\ref{r3:graphs}, use
$X=U$ as an actually observed covariate. It permits adjustment. Mapping
all $X$ to a constant loses the distinction and recovers a spurious
contrast of $0.36$. Adding an excellent predictor of a different
component of $Y$ cannot restore the discarded confounder. Conversely,
balancing representations cannot remove an unmeasured common cause
that was absent from the input. Distributional invariance across
environments needs assumptions about which mechanisms change; stable
prediction in two selected environments does not establish that an
association survives an action intervention.

\section{Consequences for the shared learning layer}
\label{r3:ml-contract}
The product should accept an identified functional or moment, an
explicit target measure, nuisance roles and a sampling model. It can then
choose compatible learners, splitting schemes and uncertainty methods.
The output must identify whether it is an ATE, an overlap-weighted
projection, a CATE curve, a complier curve or a policy value. A response
predictor can be shared across applications without transferring the
causal validity of any particular application.

Diagnostics should correspond to the mathematics: nuisance residuals
and support, score stability, fold leakage, local moment rank,
regularization sensitivity and policy-selection rules. Predictive
validation is useful for nuisance fit; it cannot certify exchangeability,
instrument exclusion, causal representation or transport to another bank.
These distinctions also determine what automated data organization is
allowed to discard or aggregate. The next chapter develops inverse
learning with endogeneity and proxies. Chapter~\ref{r3:longitudinal}
then makes the information set dynamic and shows why sequence prediction
is still only one part of the causal problem.
